\section{Introduction}
\label{sec:intro}

Style transfer on a frozen generator is treated as one task, but
users want two. Given a content prompt and a reference image, a frozen
text-to-image transformer can render the prompt's subject ``in the style of''
the reference. In practice two very different outcomes are desirable and are
conflated under one name. In one, an artist wants the reference's
\emph{scene} preserved, including its composition and arrangement, with only the
subject re-imagined in the reference's medium. This is \emph{Semantic Style
Editing}. In the other, the artist wants only the reference's \emph{look},
including its medium, texture, and palette, carried onto a freshly composed
scene, with the reference's specific layout thrown away. This is
\emph{Global Style Transfer}. The desired amount of reference-layout retention therefore depends on
the creative task.

These two outcomes motivate a common \emph{style-locality} axis.
Style locality describes how much reference layout is retained. Greater
retention supports semantic editing, while freer composition supports global
transfer. We assess this axis separately from reference-medium fidelity and
target-subject alignment, so replacing the reference subject remains an
explicit requirement in either task.

We drive this axis with one frozen-model operator. Distribution-matching
distillation separates CFG augmentation, which strengthens conditioning,
from distribution matching with a learned fake-score estimate
\citep{yin2024dmd,liu2025decoupled}. Our setting differs: we optimize a
single inference-time latent while retaining the frozen reverse process.
The native reference-conditioning contrast supplies the style direction,
and ordinary reverse steps continue the model's denoising trajectory
between latent interventions. We call this procedure \textbf{Test-Time
Score Distillation (TSD)} and its alternation \emph{reverse-process
regularization}.

TSD evaluates reference-conditioned and reference-free velocities at the
same latent state, noise level, and target text. We convert their difference
to a score direction, RMS-normalize it, and construct a detached clean-state
target for latent optimization. The target is rebuilt after each inner
update. All network weights remain frozen, with no learned style adapter,
auxiliary fake-score network, or spatial mask. Sampling resolution and
intervention settings jointly determine the style-locality operating point.

\paragraph{Contributions.}
\begin{enumerate}
\item \textbf{A common task formulation.} Semantic Style Editing and Global
Style Transfer are two goals on a style-locality axis. We serve both with
one operator and study their sampling and latent-intervention settings.
\item \textbf{Test-time score distillation with reverse-process
regularization.} We turn the frozen model's same-state reference contrast
into an RMS-normalized pseudo-target for latent optimization, interleaved
with its native reverse trajectory.
\item \textbf{A separable-axis evaluation protocol.} We evaluate style
fidelity with CSD, subject fidelity with CLIP-T, and composition transfer
with reference-layout similarity, interpreting all three together.
\item \textbf{Native-editor and dedicated-method comparisons.} We retain
the initial 40-reference study and extend it with multiple prompts, a strong
same-backbone reference-CFG control, and dedicated transfer methods. The joint readouts expose how guidance and intervention settings affect
style, subject, and composition.
\end{enumerate}

